% Gesamttest «Betragsfunktionen» — Quelle des PDFs (python3 scripts/build-lp-pdf.py betrag).
% Musterlösung und Raster: bewertungspaket.tex. Alle Werte mit python3 nachgerechnet (05.10.2026).
% Kein \lt/\gt — pdfLaTeX kennt sie nicht, < und > tun es.
\documentclass[11pt]{article}
\usepackage{../lp-druck}
\renewcommand{\lpname}{Betragsfunktionen}
\renewcommand{\lpteilgebiet}{3.6}
\renewcommand{\lpbereich}{Schwerpunktfach}
\renewcommand{\lpfuss}{Gesamttest · 25 Punkte}
\begin{document}
\lpkopf{Gesamttest}{%
  \vspace{4pt}\noindent\begin{tabularx}{\linewidth}{@{}X X@{}}
  Code (kein Name): \hrulefill & Punkte: \hrulefill\ / 25
  \end{tabularx}}

\begin{lpinfo}{Anleitung}
\textbf{Zeit:} rund 30 Minuten. \textbf{Hilfsmittel:} kein Taschenrechner; Lineal und Bleistift sind erlaubt.
Schreib zu jeder Aufgabe den \textbf{Rechenweg} auf — bewertet wird der Weg.
Danach: Lösung scannen oder fotografieren (ohne Namen und Standort) und mit dem \emph{Bewertungspaket} bewerten lassen.
\end{lpinfo}

\lpteil{Teil A · Graphen}{Kapitel 1 bis 3 · 14 Punkte}

\begin{lpaufgabe}{G1}{Die Betragsfunktion}{3 P}
\begin{enumerate}[label=(\alph*),leftmargin=*,itemsep=1pt]
\item Berechne $|3 - 10| - |-4|$.
\item Schreib $|x|$ mit Fallunterscheidung und skizziere $y = |x|$ für $-3 \le x \le 3$ (Platz unten).
\item Für welche $x$ gilt $|x| = 4.5$?
\end{enumerate}
\lpplatz{5}
\end{lpaufgabe}

\begin{lpaufgabe}{G2}{Ein Dach skizzieren}{4 P}
Skizziere $y = -2\,|x + 1| + 4$ ohne Wertetabelle. Gib den Knickpunkt und die Steigungen der beiden Äste an und berechne die Nullstellen.
\par\smallskip\centering
\begin{tikzpicture}
\begin{axis}[width=12cm,height=7.5cm,axis lines=middle,xmin=-4.5,xmax=5.5,ymin=-5.5,ymax=4.5,
  xtick={-4,-3,-2,-1,1,2,3,4,5},ytick={-5,-4,-3,-2,-1,1,2,3,4},grid=both,grid style={lplinie},minor tick num=0,axis on top,tick label style={font=\scriptsize},
  xlabel=$x$,ylabel=$y$,
  yticklabel style={fill=white,inner sep=1pt,xshift=-2pt},xticklabel style={fill=white,inner sep=1pt,yshift=-2pt}]
\end{axis}
\end{tikzpicture}
\par\raggedright
\lpplatz{3}
\end{lpaufgabe}

\begin{lpaufgabe}{G3}{Gleichung aus dem Graphen}{3 P}
Bestimme die Gleichung der abgebildeten Kurve in der Form $y = a\,|x - u| + v$. Begründe $a$ mit einem abgelesenen Punkt.
\par\smallskip\centering
\begin{tikzpicture}
\begin{axis}[width=12cm,height=6.5cm,axis lines=middle,xmin=-6.5,xmax=4.5,ymin=-2.5,ymax=3.5,
  xtick={-6,-5,-4,-3,-2,-1,1,2,3,4},ytick={-2,-1,1,2,3},grid=both,grid style={lplinie},minor tick num=0,axis on top,tick label style={font=\scriptsize},
  xlabel=$x$,ylabel=$y$,
  yticklabel style={fill=white,inner sep=1pt,xshift=-2pt},xticklabel style={fill=white,inner sep=1pt,yshift=-2pt}]
\addplot[lpblau,very thick,domain=-6.5:4.5,samples=200]{0.5*abs(x+2)-1};
\end{axis}
\end{tikzpicture}
\par\raggedright
\lpplatz{2}
\end{lpaufgabe}

\begin{lpaufgabe}{G4}{Umklappen}{4 P}
Gegeben ist $f(x) = x^2 - 2x - 3$.
\begin{enumerate}[label=(\alph*),leftmargin=*,itemsep=1pt]
\item Bestimme die Nullstellen und den Scheitel von $f$.
\lpplatz{3}
\item Skizziere $f$ und $y = |f(x)|$ in das Koordinatensystem. Gib die Knicke und den umgeklappten Scheitel an.
\end{enumerate}
\par\smallskip\centering
\begin{tikzpicture}
\begin{axis}[width=12cm,height=7.5cm,axis lines=middle,xmin=-3.5,xmax=5.5,ymin=-4.5,ymax=5.5,
  xtick={-3,-2,-1,1,2,3,4,5},ytick={-4,-3,-2,-1,1,2,3,4,5},grid=both,grid style={lplinie},minor tick num=0,axis on top,tick label style={font=\scriptsize},
  xlabel=$x$,ylabel=$y$,
  yticklabel style={fill=white,inner sep=1pt,xshift=-2pt},xticklabel style={fill=white,inner sep=1pt,yshift=-2pt}]
\end{axis}
\end{tikzpicture}
\par\raggedright
\lpplatz{2}
\end{lpaufgabe}

\lpteil{Teil B · Terme und Gleichungen}{Kapitel 4 und 5 · 11 Punkte}

\begin{lpaufgabe}{G5}{Abschnittsweise schreiben}{4 P}
\begin{enumerate}[label=(\alph*),leftmargin=*,itemsep=1pt]
\item Schreib $|2x - 5|$ ohne Betragsstriche, mit Fallunterscheidung.
\item Schreib $y = |x + 1| + |x - 4|$ abschnittsweise, mit allen drei Abschnitten. Wie hoch liegt der flache Boden?
\end{enumerate}
\lpplatz{6}
\end{lpaufgabe}

\begin{lpaufgabe}{G6}{Gleichung und Ungleichung}{4 P}
Löse rechnerisch und kontrolliere mit einer Skizze:
\begin{enumerate}[label=(\alph*),leftmargin=*,itemsep=1pt]
\item $|3x - 6| = 9$
\item $|x + 1| \ge 2$
\end{enumerate}
\lpplatz{6}
\end{lpaufgabe}

\begin{lpaufgabe}{G7}{Lösungen zählen}{3 P}
$f$ ist die Funktion aus G4. Wie viele Lösungen hat $|f(x)| = c$ für $c = 2$, $c = 4$ und $c = 5$? Begründe mit deiner Skizze aus G4 oder rechnerisch.
\lpplatz{6}
\end{lpaufgabe}
\end{document}
